
\documentclass[12pt,a4paper]{ctexart}

\usepackage[margin=2.3cm]{geometry}
\usepackage{amsmath}
\usepackage{array}
\usepackage{booktabs}
\usepackage{caption}
\usepackage{enumitem}
\usepackage{etoolbox}
\usepackage{fancyhdr}
\usepackage{float}
\usepackage{graphicx}
\usepackage{hyperref}
\usepackage{longtable}
\usepackage{tabularx}
\usepackage[table]{xcolor}
\usepackage{tikz}
\usetikzlibrary{arrows.meta,positioning,shapes.geometric,fit,calc}

\hypersetup{hidelinks}
\urlstyle{same}
\setlist{nosep}
\captionsetup{font=small}
\renewcommand{\arraystretch}{1.25}
\setlength{\parindent}{2em}
\setlength{\parskip}{0.3em}
\setlength{\tabcolsep}{4pt}
\emergencystretch=3em
\sloppy
\makeatletter
\g@addto@macro{\UrlBreaks}{\do\/\do\-\do\_\do\.\do\:}
\makeatother
\AtBeginEnvironment{longtable}{\small\setlength{\tabcolsep}{3pt}}
\AtBeginEnvironment{tabularx}{\small\setlength{\tabcolsep}{3pt}}
\AtBeginEnvironment{tabular}{\small\setlength{\tabcolsep}{3pt}}

\newcommand{\studentname}{许奕}
\newcommand{\studentid}{2351441}
\newcommand{\teamname}{许奕、王天一、马源胜、周由}
\newcommand{\coursename}{软件工程管理与经济}
\newcommand{\projectname}{基于大模型的建筑能源智能管理与运维优化系统}
\newcommand{\docversion}{V2.0}
\newcommand{\docdate}{2026年5月31日}
\newcommand{\code}[1]{{\footnotesize\nolinkurl{#1}}}
\newcolumntype{Y}{>{\raggedright\arraybackslash}X}
\newcolumntype{C}[1]{>{\centering\arraybackslash}p{#1}}
\newcolumntype{L}[1]{>{\raggedright\arraybackslash}p{#1}}

\tikzset{
  box/.style={rectangle, rounded corners, draw=cyan!55!black, fill=cyan!8, thick, minimum height=8mm, align=center, inner sep=5pt},
  process/.style={rectangle, rounded corners, draw=teal!55!black, fill=teal!8, thick, minimum height=8mm, align=center, inner sep=5pt},
  data/.style={cylinder, shape border rotate=90, aspect=0.25, draw=orange!70!black, fill=orange!10, thick, minimum height=12mm, align=center, inner sep=5pt},
  risk/.style={diamond, draw=red!60!black, fill=red!8, thick, aspect=2, align=center, inner sep=2pt},
  arrow/.style={-{Stealth[length=2.4mm]}, thick, draw=gray!70!black},
  note/.style={rectangle, draw=gray!60, fill=gray!6, rounded corners, align=left, inner sep=6pt}
}

\pagestyle{fancy}
\fancyhf{}
\fancyhead[L]{\coursename}
\fancyhead[R]{\reporttitle}
\fancyfoot[C]{\thepage}
\setlength{\headheight}{15pt}

\title{
    \vspace{-2em}
    \heiti \zihao{2} \reporttitle \\
    \vspace{1em}
    \zihao{4} 课程期末项目正式交付文档
}
\author{
    \songti \zihao{4}
    项目名称：\projectname \\
    项目组成员：\teamname \\
    负责人：\studentname \quad 学号：\studentid
}
\date{\docdate}

\newcommand{\reporttitle}{SEM 软件工程管理说明}
\begin{document}

\maketitle
\thispagestyle{fancy}

\begin{abstract}
本文档为《\projectname》的软件工程管理说明（SEM），覆盖项目范围、计划、组织分工、软件过程监控与控制、质量保证、配置管理、风险管理、沟通管理、变更管理和最终交付管理。文档从工程管理角度说明项目如何由初始骨架、两轮分工、两次整合发展为最终可运行、可测试、可提交的系统。
\end{abstract}

\tableofcontents
\newpage

\section{文档控制}
\begin{longtable}{p{4cm}p{10.5cm}}
\toprule
项目 & 内容 \\
\midrule
文档名称 & SEM 软件工程管理说明 \\
文档版本 & \docversion \\
管理范围 & 范围、计划、组织、过程控制、质量、风险、配置、沟通、交付 \\
项目组织 & 四人课程项目小组 \\
管理模式 & 文档先行、目录分工、两轮任务、集中整合、最终封版 \\
\bottomrule
\end{longtable}

\section{项目范围管理}
\subsection{本期范围}
本期范围包括：建筑能耗样例数据、数据字典、FastAPI 后端、Vue 前端、三维楼层态势、统计分析、异常解释、工单闭环、决策报告、智能问答、外部模型配置、MCP Server、自动化测试、运行脚本和最终文档。

\subsection{范围排除}
不纳入本期范围的内容包括：真实传感器接入、企业级登录权限、生产数据库、高可用部署、移动端 App、长期真实节能审计和商业合同实施。范围排除的目的不是降低质量，而是在课程周期内保证核心系统闭环可以稳定完成。

\section{组织与职责}
\begin{longtable}{p{2.5cm}p{4.2cm}p{4.2cm}p{3.4cm}}
\caption{团队分工与交付物}\\
\toprule
成员 & 主要职责 & 核心交付物 & 集成边界 \\
\midrule
\endfirsthead
\toprule
成员 & 主要职责 & 核心交付物 & 集成边界 \\
\midrule
\endhead
许奕 & 项目规划、整合、测试、文档、最终封版 & README、任务说明、整合总结、最终文档、测试验证 & 统一接口、统一文档、最终质量控制 \\
王天一 & 后端与架构 & REST API、分析服务、MCP Server、后端测试 & 保持 API 契约稳定 \\
马源胜 & 数据与 AI & 样例数据、数据源调研、知识库、模型配置 & 数据字段与知识库口径稳定 \\
周由 & 前端与集成 & Vue 页面、图表、3D 风险、前后端联调 & 通过 \code{api.js} 调用接口 \\
\bottomrule
\end{longtable}

\subsection{RACI 职责矩阵}
\begin{longtable}{L{3.4cm}C{1.7cm}C{1.7cm}C{1.7cm}C{1.7cm}L{3.2cm}}
\caption{RACI 职责矩阵}\\
\toprule
活动 & 许奕 & 王天一 & 马源胜 & 周由 & 说明 \\
\midrule
\endfirsthead
\toprule
活动 & 许奕 & 王天一 & 马源胜 & 周由 & 说明 \\
\midrule
\endhead
需求与范围 & A/R & C & C & C & 负责人统一范围和验收口径 \\
后端接口 & C & A/R & C & C & 后端负责实现与测试 \\
数据与知识库 & C & C & A/R & C & 数据负责人保证字段和知识库质量 \\
前端页面 & C & C & C & A/R & 前端负责页面和交互 \\
系统整合 & A/R & C & C & C & 集中整合减少冲突 \\
测试验收 & A/R & R & C & R & 自动化与人工验收结合 \\
最终文档 & A/R & C & C & C & 正式文档由负责人统一封版 \\
\bottomrule
\end{longtable}

其中 A 表示最终负责，R 表示执行负责，C 表示参与咨询。该矩阵的作用是把“谁最终负责”和“谁具体执行”分开，避免多人协作中出现职责空白。

\section{项目计划}
\subsection{阶段划分}
\begin{figure}[H]
\centering
\resizebox{0.96\textwidth}{!}{%
\begin{tikzpicture}[node distance=7mm and 6mm]
\node[process] (init) {初始化\\骨架与 README};
\node[process, right=of init] (r1) {第一次任务\\三方独立开发};
\node[process, right=of r1] (i1) {第一次整合\\补齐基础闭环};
\node[process, right=of i1] (r2) {第二次任务\\强化演示能力};
\node[process, below=of r2] (i2) {第二次整合\\系统主线完成};
\node[process, left=of i2] (polish) {封版优化\\3D/工单/LLM/MCP};
\node[process, left=of polish] (docs) {最终文档\\SRS/SDD/SEE/SEM};
\node[process, left=of docs] (submit) {提交准备\\PDF/源码/演示};
\draw[arrow] (init) -- (r1);
\draw[arrow] (r1) -- (i1);
\draw[arrow] (i1) -- (r2);
\draw[arrow] (r2) -- (i2);
\draw[arrow] (i2) -- (polish);
\draw[arrow] (polish) -- (docs);
\draw[arrow] (docs) -- (submit);
\end{tikzpicture}
}
\caption{项目阶段路径}
\end{figure}

\subsection{里程碑}
\begin{longtable}{p{3cm}p{5.4cm}p{5.5cm}}
\caption{主要里程碑}\\
\toprule
里程碑 & 完成标准 & 证据 \\
\midrule
\endfirsthead
\toprule
里程碑 & 完成标准 & 证据 \\
\midrule
\endhead
M1 仓库初始化 & 目录结构、README、环境模板、基础文档完成 & 根目录、\code{docs/}、\code{.env.example} \\
M2 第一次整合 & 后端、前端、数据 AI 初步连通 & \code{docs/11-first-integration-summary.md} \\
M3 第二次整合 & 主要页面和接口可运行，功能闭环形成 & \code{docs/12-second-integration-summary.md} \\
M4 封版优化 & 数据扩展、3D 联动、工单持久化、统计分析优化 & 前端页面和后端测试 \\
M5 文档封版 & SRS、SDD、SEE、SEM 等正式文档生成 PDF & \code{docs/final-latex/pdf/} \\
M6 最终提交 & 源码、文档、演示材料按要求整理 & \code{期末提交材料/} \\
\bottomrule
\end{longtable}

\section{软件过程监控与控制}
\subsection{Git 控制}
项目使用 GitHub 作为统一仓库。关键控制规则如下：真实 \code{.env} 不提交；每次整合前先拉取远程；整合后运行检查脚本；功能稳定后再提交；提交信息应能表达变化内容；最终提交材料排除缓存、依赖目录和个人无关作业文件。

\subsection{接口控制}
REST API 契约以 \code{docs/06-api-contract.md} 为基础，前端只通过 \code{frontend/src/lib/api.js} 调用接口，MCP Server 复用后端服务层。这样可以降低多人协作中“前端字段名”和“后端字段名”不一致的风险。

\subsection{自动化质量控制}
项目提供 \code{scripts/check-project.ps1} 作为质量验证入口，覆盖 Python 语法、后端 pytest 和前端构建。脚本失败即退出，避免出现测试失败但仍显示成功的误判。

\subsection{过程监控闭环}
\begin{figure}[H]
\centering
\resizebox{0.86\textwidth}{!}{%
\begin{tikzpicture}[x=1cm,y=1cm,every node/.style={font=\small}]
\node[process, minimum width=2.7cm, minimum height=0.9cm] (plan) at (0,1.6) {1 计划\\任务边界};
\node[process, minimum width=2.7cm, minimum height=0.9cm] (do) at (3.8,1.6) {2 执行\\分模块开发};
\node[process, minimum width=2.7cm, minimum height=0.9cm] (check) at (7.6,1.6) {3 检查\\测试与评审};
\node[process, minimum width=2.7cm, minimum height=0.9cm] (act) at (7.6,-0.4) {4 纠偏\\修复与整合};
\node[process, minimum width=2.7cm, minimum height=0.9cm] (base) at (3.8,-0.4) {5 基线\\提交与文档};
\draw[arrow] (plan) -- (do);
\draw[arrow] (do) -- (check);
\draw[arrow] (check) -- (act);
\draw[arrow] (act) -- (base);
\draw[arrow] (base.west) -- ++(-3.8,0) -- (plan.south);
\end{tikzpicture}
}
\caption{软件过程监控与控制闭环}
\end{figure}
每轮整合均以测试通过、文档同步和安全检查作为进入下一轮的条件。该闭环强调先冻结任务边界，再执行开发、测试评审、修复整合和基线归档，避免后期反复增加范围导致提交质量失控。

\section{质量管理}
\subsection{质量目标}
\begin{enumerate}
    \item 系统能够按 README 在普通 Windows 环境复现运行。
    \item 后端接口测试、服务测试、MCP 测试和 LLM 封装测试通过。
    \item 前端生产构建通过，核心页面无阻断错误。
    \item 系统展示路径覆盖项目亮点，并能应对刷新、筛选、点击、导出和问答。
    \item 最终文档使用正式口吻，面向教师评审，不包含对话式说明。
    \item 最终材料不泄露真实 API Key。
\end{enumerate}

\subsection{质量保证矩阵}
\begin{longtable}{p{3.2cm}p{4.8cm}p{5.5cm}}
\caption{质量保证矩阵}\\
\toprule
质量活动 & 控制内容 & 输出证据 \\
\midrule
\endfirsthead
\toprule
质量活动 & 控制内容 & 输出证据 \\
\midrule
\endhead
需求确认 & 功能是否覆盖课程要求和系统展示路径 & SRS、验收报告 \\
设计确认 & 架构是否清晰、接口是否一致、数据是否可追溯 & SDD、API 文档 \\
代码验证 & 路由、服务、前端页面、脚本是否能运行 & 源码目录、构建结果 \\
自动化测试 & 后端接口、服务、MCP、LLM 封装 & pytest 结果 \\
人工验收 & 浏览器操作、3D 联动、工单持久化、报告生成 & 测试与验收报告 \\
安全控制 & \code{.env}、API Key、运行期数据是否被排除 & \code{.gitignore}、交付材料 \\
\bottomrule
\end{longtable}

\section{风险管理}
\begin{longtable}{p{3.2cm}p{1.6cm}p{1.6cm}p{5.2cm}p{2.3cm}}
\caption{风险登记表}\\
\toprule
风险 & 概率 & 影响 & 应对措施 & 当前状态 \\
\midrule
\endfirsthead
\toprule
风险 & 概率 & 影响 & 应对措施 & 当前状态 \\
\midrule
\endhead
成员提交质量不一致 & 中 & 高 & 分工文件写清边界，整合阶段统一补齐 & 已控制 \\
多人合并冲突 & 中 & 中 & 按目录和模块分工，减少同时改同一文件 & 已控制 \\
数据量不足导致演示单薄 & 中 & 高 & 扩展至 2,976 条、3 月至 6 月、派生楼层设备维度 & 已控制 \\
系统像写死数据 & 中 & 高 & 增加筛选、楼层分析、工单持久化、报告和 MCP 调用 & 已控制 \\
外部模型 Key 泄露 & 中 & 高 & 真实 Key 仅放本地 \code{.env}，发布前进行模式扫描 & 已控制 \\
外部模型不可用 & 中 & 中 & 本地知识库兜底，LLM 可配置关闭 & 已控制 \\
MCP 依赖影响运行 & 低 & 中 & requirements 声明依赖，提供单独启动脚本和测试 & 已控制 \\
PPT/视频材料衔接不足 & 中 & 中 & 软件侧交付内容完整，展示材料按最终系统截图和文档组织 & 可控制 \\
\bottomrule
\end{longtable}

\section{配置管理}
\subsection{配置项}
\begin{longtable}{p{3.5cm}p{5cm}p{5.2cm}}
\caption{配置项说明}\\
\toprule
类别 & 配置项 & 管理规则 \\
\midrule
\endfirsthead
\toprule
类别 & 配置项 & 管理规则 \\
\midrule
\endhead
源码 & \code{backend/}、\code{frontend/}、\code{scripts/} & 纳入 Git 管理，版本发布前运行验证脚本 \\
数据 & \code{data/samples/}、\code{data/dictionaries/} & 样例数据与字典纳入仓库，运行期数据排除 \\
知识库 & \code{knowledge_base/} & 作为问答和 RAG 素材纳入仓库 \\
环境模板 & \code{.env.example} & 只存占位变量，不含真实密钥 \\
真实环境 & \code{.env} & 本地私有文件，不提交、不打包 \\
运行状态 & \code{data/runtime/} & 本地生成，不提交 \\
最终材料 & \code{期末提交材料/} & 统一整理源码、PDF、说明和展示材料 \\
\bottomrule
\end{longtable}

\section{变更管理}
项目后期采用集中优化方式。原因是系统主链路已经形成，继续分散开发会增加冲突和回归风险。后期变更原则为：只接受能明显提升系统展示质量、文档质量或验收通过率的变更；每次变更后确认是否影响 README、API、测试和交付材料；涉及密钥、数据和依赖的变更必须额外确认安全性。

\section{交付管理}
\subsection{最终交付物}
最终交付应包含：源码、SRS、SDD/SDS、SEE、SEM、测试与验收报告、用户手册与部署说明、数据与接口说明、最终提交说明、PPT、演示视频或视频链接。

\subsection{排除项}
最终压缩包必须排除：\code{.env}、\code{frontend/node_modules/}、\code{frontend/dist/}、\code{data/runtime/}、\code{.pytest_cache/}、个人无关实验文件、临时日志和真实 API Key。

\section{软件过程改进}

本项目采用两轮开发加最终整合的迭代演进模式。第一轮完成后进行了系统的过程回顾，识别出多个需要改进的领域，并在第二轮中进行了针对性改进。

\subsection{第一轮过程回顾与问题识别}

第一次整合完成后，项目对三位成员的独立交付物进行了质量审计，识别出以下过程问题：

\begin{longtable}{L{2.8cm}L{4.2cm}L{6.8cm}}
\caption{第一轮过程问题与改进需求}\\
\toprule
领域 & 识别的问题 & 对系统的影响 \\
\midrule
\endfirsthead
\toprule
领域 & 识别的问题 & 对系统的影响 \\
\midrule
\endhead
数据规模 & 样例数据仅覆盖单日范围，难以支撑趋势展示和多时段异常分析。 & 统计图表变化不明显，演示说服力不足。 \\
界面表达 & 前端部分交互按钮标记为“计划中”但无实际功能实现；存在“文档大于实现”的问题。 & 评审人员打开页面后实际可操作功能与说明不一致。 \\
问答能力 & 后端问答接口为纯占位，无法基于项目数据和知识库给出有意义的回答。 & 智能问答页面缺乏演示价值。 \\
交付纪律 & 个别成员意外修改了不属于自己负责范围的文件（如 \code{.gitignore}、根目录文件）。 & 增加整合排查成本，造成仓库基线短暂不一致。 \\
\bottomrule
\end{longtable}

\subsection{第二轮过程改进措施}

基于第一轮的回顾发现，项目在第二轮中实施了以下过程改进。改进内容基于第二次任务统一说明中对各成员的具体要求及第二次整合总结中的验证结果：

\begin{longtable}{L{2.4cm}L{3.2cm}L{4.0cm}L{4.4cm}}
\caption{第二轮过程改进要点}\\
\toprule
改进项 & 改进前状态 & 改进措施 & 改进后效果 \\
\midrule
\endfirsthead
\toprule
改进项 & 改进前状态 & 改进措施 & 改进后效果 \\
\midrule
\endhead
数据扩充 & 单日样例数据 & 将数据集扩展至 2026-03-01 至 2026-06-01，覆盖多建筑、多楼层、多时段，最终形成 2,976 条记录。 & 趋势图有明显变化，异常分析有足够样本量支撑多维演示。 \\
图表增强 & 前端图表为简单占位 & 接入正式图表组件，实现趋势折线图、建筑对比柱状图、COP 排名和异常原因统计图；并改为按需注册 ECharts、统一管理实例生命周期。 & 统计分析页图表可随筛选范围动态刷新，展示效果接近正式产品。 \\
问答升级 & 问答接口为占位 & 新增知识检索服务 \code{knowledge\_search\_service.py}，将知识库 Markdown 文件接入回答引用逻辑。 & 无外部模型时也能给出基于项目知识库的有意义回答。 \\
建筑分析 & 分析维度较粗 & 新增建筑深度分析服务与接口，补充建筑画像和异常场景参考数据。 & 分析维度从建筑级细化到楼层和设备级。 \\
真实交互 & 部分按钮无功能 & 补齐建筑筛选、时间范围筛选、楼层筛选和 CSV 导出功能，前端围绕已稳定 API 实现。 & 演示路径内所有关键操作均可完成，不出现“计划中”的空按钮。 \\
交付纪律 & 偶有越界修改 & 第二轮任务说明中明确限定每人只能修改自己负责的目录，禁止交叉写文件；前端不新建占位功能。 & 第二轮整合时未再发现越界修改问题。 \\
接口稳定性 & 第一轮接口已稳定 & 第二轮要求“可以新增接口但不许改坏旧接口”，后端新增知识检索和建筑深度分析接口，保留原有接口不变。 & 第一轮已接通的页面在第二轮和最终版本中未出现回退。 \\
\bottomrule
\end{longtable}

\subsection{过程改进度量}

为量化过程改进效果，项目选取了若干可度量指标进行前后对比。数据来源为第一次整合总结（\code{docs/11-}）、第二次整合总结（\code{docs/12-}）及最终封版后的验证结果：

\begin{longtable}{L{3.4cm}L{2.2cm}L{2.2cm}L{6.2cm}}
\caption{过程改进度量指标}\\
\toprule
度量指标 & 第一轮 & 第二轮/最终 & 改进说明 \\
\midrule
\endfirsthead
\toprule
度量指标 & 第一轮 & 第二轮/最终 & 改进说明 \\
\midrule
\endhead
后端测试通过数 & 33 & 75 & 第一次整合 33 passed → 第二次整合 47 passed → 最终（含 MCP 测试）75 passed \\
前端构建 & 通过 & 通过 & 第二轮修复了 \code{package-lock.json} 未同步导致的构建失败，且对图表组件做了按需加载优化 \\
后端接口数 & 约 12 个 & 约 21 个 & 新增知识检索、建筑深度分析、异常解释、工单、报告、MCP 等接口 \\
演示路径节点 & 约 4 步 & 约 8 步 & 从总览、查询、分析到问答形成完整闭环，增加 3D 联动和工单节点 \\
越界文件修改 & 2 处 & 0 处 & 第二轮任务边界约束有效消除了越界问题 \\
\bottomrule
\end{longtable}

\subsection{过程改进的经验教训}

\begin{enumerate}
    \item \textbf{明确的目录边界是多人协作的基础}：通过限定每人只修改自己负责的目录，显著降低了合并冲突和意外破坏已有功能的风险。
    \item \textbf{先稳定接口再增强功能}：第一轮优先保证前后端接口契约稳定，第二轮在此基础之上扩展分析深度和交互体验，避免了功能扩展导致已有页面回退。
    \item \textbf{用测试脚本作为质量闸门}：每次整合后运行 \code{check-project.ps1} 一键验证，确保新提交不会破坏已有测试。后端测试从 33 项逐步增长到 75 项，每次增量都有明确覆盖范围。
    \item \textbf{及早发现文档与实现不一致}：第一轮暴露的“文档大于实现”问题促使第二轮严格执行“只写真实实现、不写计划中功能”的原则。
    \item \textbf{数据规模是演示效果的基础}：从单日数据扩展到三个月 2,976 条记录是提升系统可信度最关键的一步，应尽早规划而非留到最终封版。
\end{enumerate}

\section{项目收尾}

\subsection{交付物确认}

项目已按教师课程要求完成全部交付物的整理和确认：

\begin{longtable}{L{3.0cm}L{5.0cm}L{5.5cm}}
\caption{最终交付物确认清单}\\
\toprule
交付物类别 & 交付物明细 & 确认状态 \\
\midrule
\endfirsthead
\toprule
交付物类别 & 交付物明细 & 确认状态 \\
\midrule
\endhead
设计规范 & SRS 软件需求规格说明书、SDD/SDS 软件设计说明书 & 已生成 PDF \\
经济与管理 & SEE 软件经济分析与评价、SEM 软件工程管理说明 & 已生成 PDF \\
测试与验收 & 测试与验收报告、后端 75 项测试通过 & 已确认 \\
用户与部署 & 用户手册与部署说明、README 运行说明 & 已确认 \\
数据与接口 & 数据接口与 MCP 说明、样例数据与数据字典 & 已确认 \\
提交说明 & 最终提交说明、提交材料目录 & 已确认 \\
源码 & \code{backend/}、\code{frontend/}、\code{data/}、\code{scripts/}、\code{knowledge\_base/} & 已确认可运行 \\
演示材料 & PPT、演示视频 & 以最终系统为准 \\
\bottomrule
\end{longtable}

\subsection{未完成事项与后续建议}

本期课程项目范围内的工作已经完成。对于项目后续可能的延续方向，建议关注：
\begin{enumerate}
    \item 将 CSV 数据存储替换为 SQLite 或 PostgreSQL，为更大数据量做准备。
    \item 增加用户登录和基于角色的权限控制，支持真实部署环境。
    \item 接入真实建筑传感器数据源，减少对模拟数据的依赖。
    \item 将可解释规则替换或增强为机器学习模型，提升异常检测智能水平。
    \item 完善多租户 SaaS 体系，为商业化运营做准备。
\end{enumerate}

\subsection{项目总结}

从软件工程管理角度看，本项目完成了从范围定义、任务分解、两轮多人开发、两次集中整合、过程监控、质量控制、风险收敛到最终交付的完整工程闭环。项目在开发过程中识别了协作问题（越界修改、文档与实现不一致、数据量不足、问答占位等），并在第二轮中有针对性地进行了过程改进，改进效果通过测试数量增长、接口数量增长和演示路径扩展等可度量指标得到了验证。

系统实现已形成可运行的 Web 系统，正式文档已形成完整 PDF 交付包，交付材料已按教师要求整理。项目具备课程期末提交条件。

\clearpage
\section{补充工程管理说明}
\subsection{阶段计划与控制点}
项目采用两轮任务分配加最终整合的方式推进。管理重点不是追求形式上的每日进度，而是保证每个阶段都有明确输入、输出、检查项和纠偏动作。
\begin{longtable}{L{2.0cm}L{3.2cm}L{5.2cm}L{4.2cm}}
\caption{阶段计划、产出与控制点}\\
\toprule
阶段 & 目标 & 主要产出 & 控制点 \\
\midrule
\endfirsthead
\toprule
阶段 & 目标 & 主要产出 & 控制点 \\
\midrule
\endhead
S1 & 项目骨架 & README、目录结构、环境示例、基础文档和样例数据。 & 能够说明项目要做什么。 \\
S2 & 第一次分工 & 后端、数据 AI、前端三个方向的独立任务说明。 & 边界清楚，减少合并冲突。 \\
S3 & 第一次整合 & 统一接口、补齐测试、修复质量不足、形成第一版系统。 & 运行脚本通过。 \\
S4 & 第二次分工 & 围绕系统完善、AI、3D、工单、报告和文档继续细化。 & 不再大规模改架构。 \\
S5 & 第二次整合 & 合并成员提交，统一代码风格，补充缺失逻辑。 & 后端测试和前端构建通过。 \\
S6 & 最终优化 & 扩充数据、优化 3D、工单持久化、统计分析和文档。 & 演示路径完整。 \\
S7 & 正式文档 & SRS、SDD、SEE、SEM、测试、用户手册、数据接口和提交说明。 & PDF 可读且无密钥。 \\
S8 & 最终交付 & 提交材料目录、源码副本、PDF、PPT 和视频。 & 按教师要求打包。 \\
\bottomrule
\end{longtable}

\subsection{风险登记册}
风险管理关注对最终提交影响最大的事项，包括运行风险、文档风险、密钥风险和协作风险。
\begin{longtable}{L{1.6cm}L{3.0cm}L{4.8cm}L{4.8cm}}
\caption{风险登记册}\\
\toprule
编号 & 风险 & 影响 & 应对措施 \\
\midrule
\endfirsthead
\toprule
编号 & 风险 & 影响 & 应对措施 \\
\midrule
\endhead
R-01 & 成员提交质量不均 & 接口、数据和页面可能无法直接整合。 & 整合者补齐缺口并用测试脚本作为质量闸门。 \\
R-02 & 合并冲突 & 多人同时改同一文件导致返工。 & 任务边界按文件划分，尽量减少交叉修改。 \\
R-03 & 数据量过少 & 演示显得像写死数据。 & 扩展到多建筑、多楼层、多时段。 \\
R-04 & 业务逻辑单薄 & 系统容易退化为单纯查询页面。 & 增加楼层分析、设备分析、异常解释、工单闭环和报告。 \\
R-05 & 外部模型不可用 & 问答功能演示失败。 & 本地知识库优先，外部模型可选。 \\
R-06 & 真实密钥泄露 & GitHub 公开后造成安全和费用风险。 & .env 不入库，发布前模式搜索。 \\
R-07 & 文档版面错误 & PDF 图表重叠或表格出页影响评分。 & 使用 LaTeX 重写并检查 Overfull 日志和页面预览。 \\
R-08 & 演示刷新丢工单 & 交互逻辑显得不完整。 & 工单写入运行期 JSON。 \\
R-09 & 前端构建失败 & 最终源码不可复现。 & 发布前运行 Vite build。 \\
R-10 & 后端测试失败 & 接口口径不可信。 & pytest 作为发布前必须验证项。 \\
R-11 & PPT 与系统不一致 & 系统展示和实际功能脱节。 & PPT 以最终系统截图和文档为准。 \\
R-12 & 压缩包结构混乱 & 评审人员难以找到正式文档和源码。 & 使用统一提交材料目录。 \\
R-13 & 运行说明不足 & 拉取代码后无法启动。 & README 和用户手册写明完整步骤。 \\
R-14 & 系统过度依赖单机状态 & 换机器运行时工单数据缺失。 & 说明运行期数据可自动生成。 \\
R-15 & 图表含义不清 & 评审无法理解业务价值。 & 文档为每张图补充文字解释。 \\
R-16 & 时间线不真实 & 课程过程材料可能与实际开发不完全一致。 & 最终文档聚焦工程成果和可验收内容。 \\
\bottomrule
\end{longtable}

\subsection{质量检查矩阵}
\begin{longtable}{L{2.4cm}L{4.0cm}L{4.5cm}L{3.5cm}}
\caption{质量检查矩阵}\\
\toprule
检查对象 & 检查内容 & 通过标准 & 证据 \\
\midrule
\endfirsthead
\toprule
检查对象 & 检查内容 & 通过标准 & 证据 \\
\midrule
\endhead
README & 环境、启动、测试、密钥说明 & 同学按说明可运行系统。 & 人工复现 \\
后端 & 语法、接口、服务测试 & pytest 全部通过。 & 75 passed \\
前端 & 构建、页面加载、交互 & Vite build 成功，页面可演示。 & 构建日志 \\
数据 & 规模、字段、时间范围 & 覆盖 2026-03-01 至 2026-06-01。 & 数据字典 \\
3D & 旋转、颜色、点击联动 & 异常楼层跳转统计分析。 & 浏览器验收 \\
工单 & 创建、完成、排序、持久化 & 刷新后状态保留。 & 接口和人工测试 \\
AI 问答 & 本地回答、模型选择、回退 & 无 Key 时不阻断演示。 & 问答测试 \\
MCP & 工具、资源、提示词 & 可通过 MCP 调用项目能力。 & MCP 测试 \\
SRS & 需求、用例、追踪矩阵 & 可说明系统范围。 & PDF \\
SDD & 架构、模块、数据、状态机 & 可说明实现方案。 & PDF \\
SEE & 成本、收益、定价、评价 & 覆盖教师经济分析要求。 & PDF \\
SEM & 范围、过程、风险、质量 & 覆盖工程管理要求。 & PDF \\
安全 & .env 和 Key & 最终材料不包含真实密钥。 & 模式搜索 \\
交付包 & 源码、文档、PPT、视频 & 结构清晰，可直接上传。 & 目录说明 \\
\bottomrule
\end{longtable}

\subsection{管理控制项}
以下控制项覆盖项目过程监控和最终交付管理，是从“能做出来”到“能稳定交付”的管理保证。
\begin{longtable}{L{1.8cm}L{3.2cm}L{6.4cm}L{4.0cm}}
\caption{管理控制项}\\
\toprule
编号 & 控制项 & 控制说明 & 责任输出 \\
\midrule
\endfirsthead
\toprule
编号 & 控制项 & 控制说明 & 责任输出 \\
\midrule
\endhead
MC-01 & 任务分解 & 对“任务分解”设置明确控制方式，发现问题后由项目整合人员修复或协调相关成员补齐。 & 过程记录、提交记录或最终文档。 \\
MC-02 & 成员边界 & 对“成员边界”设置明确控制方式，发现问题后由项目整合人员修复或协调相关成员补齐。 & 过程记录、提交记录或最终文档。 \\
MC-03 & 第一次整合 & 对“第一次整合”设置明确控制方式，发现问题后由项目整合人员修复或协调相关成员补齐。 & 过程记录、提交记录或最终文档。 \\
MC-04 & 第二次整合 & 对“第二次整合”设置明确控制方式，发现问题后由项目整合人员修复或协调相关成员补齐。 & 过程记录、提交记录或最终文档。 \\
MC-05 & 最终优化 & 对“最终优化”设置明确控制方式，发现问题后由项目整合人员修复或协调相关成员补齐。 & 过程记录、提交记录或最终文档。 \\
MC-06 & 接口控制 & 对“接口控制”设置明确控制方式，发现问题后由项目整合人员修复或协调相关成员补齐。 & 过程记录、提交记录或最终文档。 \\
MC-07 & 数据控制 & 对“数据控制”设置明确控制方式，发现问题后由项目整合人员修复或协调相关成员补齐。 & 过程记录、提交记录或最终文档。 \\
MC-08 & 代码验证 & 对“代码验证”设置明确控制方式，发现问题后由项目整合人员修复或协调相关成员补齐。 & 过程记录、提交记录或最终文档。 \\
MC-09 & 测试闸门 & 对“测试闸门”设置明确控制方式，发现问题后由项目整合人员修复或协调相关成员补齐。 & 过程记录、提交记录或最终文档。 \\
MC-10 & 文档闸门 & 对“文档闸门”设置明确控制方式，发现问题后由项目整合人员修复或协调相关成员补齐。 & 过程记录、提交记录或最终文档。 \\
MC-11 & 安全闸门 & 对“安全闸门”设置明确控制方式，发现问题后由项目整合人员修复或协调相关成员补齐。 & 过程记录、提交记录或最终文档。 \\
MC-12 & 交付闸门 & 对“交付闸门”设置明确控制方式，发现问题后由项目整合人员修复或协调相关成员补齐。 & 过程记录、提交记录或最终文档。 \\
MC-13 & 风险跟踪 & 对“风险跟踪”设置明确控制方式，发现问题后由项目整合人员修复或协调相关成员补齐。 & 过程记录、提交记录或最终文档。 \\
MC-14 & 变更记录 & 对“变更记录”设置明确控制方式，发现问题后由项目整合人员修复或协调相关成员补齐。 & 过程记录、提交记录或最终文档。 \\
MC-15 & 沟通机制 & 对“沟通机制”设置明确控制方式，发现问题后由项目整合人员修复或协调相关成员补齐。 & 过程记录、提交记录或最终文档。 \\
MC-16 & 质量目标 & 对“质量目标”设置明确控制方式，发现问题后由项目整合人员修复或协调相关成员补齐。 & 过程记录、提交记录或最终文档。 \\
MC-17 & 配置管理 & 对“配置管理”设置明确控制方式，发现问题后由项目整合人员修复或协调相关成员补齐。 & 过程记录、提交记录或最终文档。 \\
MC-18 & 版本管理 & 对“版本管理”设置明确控制方式，发现问题后由项目整合人员修复或协调相关成员补齐。 & 过程记录、提交记录或最终文档。 \\
MC-19 & 依赖管理 & 对“依赖管理”设置明确控制方式，发现问题后由项目整合人员修复或协调相关成员补齐。 & 过程记录、提交记录或最终文档。 \\
MC-20 & 运行环境 & 对“运行环境”设置明确控制方式，发现问题后由项目整合人员修复或协调相关成员补齐。 & 过程记录、提交记录或最终文档。 \\
MC-21 & 系统展示 & 对“系统展示”设置明确控制方式，发现问题后由项目整合人员修复或协调相关成员补齐。 & 过程记录、提交记录或最终文档。 \\
MC-22 & PPT 协同 & 对“PPT 协同”设置明确控制方式，发现问题后由项目整合人员修复或协调相关成员补齐。 & 过程记录、提交记录或最终文档。 \\
MC-23 & 视频协同 & 对“视频协同”设置明确控制方式，发现问题后由项目整合人员修复或协调相关成员补齐。 & 过程记录、提交记录或最终文档。 \\
MC-24 & 最终打包 & 对“最终打包”设置明确控制方式，发现问题后由项目整合人员修复或协调相关成员补齐。 & 过程记录、提交记录或最终文档。 \\
MC-25 & 密钥排查 & 对“密钥排查”设置明确控制方式，发现问题后由项目整合人员修复或协调相关成员补齐。 & 过程记录、提交记录或最终文档。 \\
MC-26 & PDF 版式 & 对“PDF 版式”设置明确控制方式，发现问题后由项目整合人员修复或协调相关成员补齐。 & 过程记录、提交记录或最终文档。 \\
MC-27 & 图表版式 & 对“图表版式”设置明确控制方式，发现问题后由项目整合人员修复或协调相关成员补齐。 & 过程记录、提交记录或最终文档。 \\
MC-28 & 表格版式 & 对“表格版式”设置明确控制方式，发现问题后由项目整合人员修复或协调相关成员补齐。 & 过程记录、提交记录或最终文档。 \\
MC-29 & README 说明 & 对“README 说明”设置明确控制方式，发现问题后由项目整合人员修复或协调相关成员补齐。 & 过程记录、提交记录或最终文档。 \\
MC-30 & 部署说明 & 对“部署说明”设置明确控制方式，发现问题后由项目整合人员修复或协调相关成员补齐。 & 过程记录、提交记录或最终文档。 \\
MC-31 & 验收路径 & 对“验收路径”设置明确控制方式，发现问题后由项目整合人员修复或协调相关成员补齐。 & 过程记录、提交记录或最终文档。 \\
MC-32 & 问题记录 & 对“问题记录”设置明确控制方式，发现问题后由项目整合人员修复或协调相关成员补齐。 & 过程记录、提交记录或最终文档。 \\
MC-33 & 纠偏动作 & 对“纠偏动作”设置明确控制方式，发现问题后由项目整合人员修复或协调相关成员补齐。 & 过程记录、提交记录或最终文档。 \\
MC-34 & 经验总结 & 对“经验总结”设置明确控制方式，发现问题后由项目整合人员修复或协调相关成员补齐。 & 过程记录、提交记录或最终文档。 \\
MC-35 & 后续维护 & 对“后续维护”设置明确控制方式，发现问题后由项目整合人员修复或协调相关成员补齐。 & 过程记录、提交记录或最终文档。 \\
\bottomrule
\end{longtable}

\subsection{过程证据说明}
过程证据说明用于体现项目具备需求、设计、开发、测试、经济分析、管理控制和交付整理的完整过程。
\begin{longtable}{L{1.7cm}L{4.0cm}L{6.2cm}L{3.6cm}}
\caption{过程证据说明}\\
\toprule
编号 & 证据项 & 说明 & 存放位置 \\
\midrule
\endfirsthead
\toprule
编号 & 证据项 & 说明 & 存放位置 \\
\midrule
\endhead
PE-01 & 第一次任务说明 & “第一次任务说明”用于支撑软件工程管理过程，能够说明项目从任务分解到最终交付的演进路径。 & 仓库文档或交付材料目录。 \\
PE-02 & 第二次任务说明 & “第二次任务说明”用于支撑软件工程管理过程，能够说明项目从任务分解到最终交付的演进路径。 & 仓库文档或交付材料目录。 \\
PE-03 & 整合说明 & “整合说明”用于支撑软件工程管理过程，能够说明项目从任务分解到最终交付的演进路径。 & 仓库文档或交付材料目录。 \\
PE-04 & 接口契约 & “接口契约”用于支撑软件工程管理过程，能够说明项目从任务分解到最终交付的演进路径。 & 仓库文档或交付材料目录。 \\
PE-05 & 数据字典 & “数据字典”用于支撑软件工程管理过程，能够说明项目从任务分解到最终交付的演进路径。 & 仓库文档或交付材料目录。 \\
PE-06 & 测试计划 & “测试计划”用于支撑软件工程管理过程，能够说明项目从任务分解到最终交付的演进路径。 & 仓库文档或交付材料目录。 \\
PE-07 & 验收报告 & “验收报告”用于支撑软件工程管理过程，能够说明项目从任务分解到最终交付的演进路径。 & 仓库文档或交付材料目录。 \\
PE-08 & 运行脚本 & “运行脚本”用于支撑软件工程管理过程，能够说明项目从任务分解到最终交付的演进路径。 & 仓库文档或交付材料目录。 \\
PE-09 & 环境示例 & “环境示例”用于支撑软件工程管理过程，能够说明项目从任务分解到最终交付的演进路径。 & 仓库文档或交付材料目录。 \\
PE-10 & Git 忽略规则 & “Git 忽略规则”用于支撑软件工程管理过程，能够说明项目从任务分解到最终交付的演进路径。 & 仓库文档或交付材料目录。 \\
PE-11 & 安全控制 & “安全控制”用于支撑软件工程管理过程，能够说明项目从任务分解到最终交付的演进路径。 & 仓库文档或交付材料目录。 \\
PE-12 & 构建日志 & “构建日志”用于支撑软件工程管理过程，能够说明项目从任务分解到最终交付的演进路径。 & 仓库文档或交付材料目录。 \\
PE-13 & 测试日志 & “测试日志”用于支撑软件工程管理过程，能够说明项目从任务分解到最终交付的演进路径。 & 仓库文档或交付材料目录。 \\
PE-14 & 正式 PDF & “正式 PDF”用于支撑软件工程管理过程，能够说明项目从任务分解到最终交付的演进路径。 & 仓库文档或交付材料目录。 \\
PE-15 & LaTeX 源文件 & “LaTeX 源文件”用于支撑软件工程管理过程，能够说明项目从任务分解到最终交付的演进路径。 & 仓库文档或交付材料目录。 \\
PE-16 & 最终交付说明 & “最终交付说明”用于支撑软件工程管理过程，能够说明项目从任务分解到最终交付的演进路径。 & 仓库文档或交付材料目录。 \\
PE-17 & 系统展示脚本 & “系统展示脚本”用于支撑软件工程管理过程，能够说明项目从任务分解到最终交付的演进路径。 & 仓库文档或交付材料目录。 \\
PE-18 & PPT 材料 & “PPT 材料”用于支撑软件工程管理过程，能够说明项目从任务分解到最终交付的演进路径。 & 仓库文档或交付材料目录。 \\
PE-19 & 视频材料 & “视频材料”用于支撑软件工程管理过程，能够说明项目从任务分解到最终交付的演进路径。 & 仓库文档或交付材料目录。 \\
PE-20 & 交付目录 & “交付目录”用于支撑软件工程管理过程，能够说明项目从任务分解到最终交付的演进路径。 & 仓库文档或交付材料目录。 \\
PE-21 & 源码副本 & “源码副本”用于支撑软件工程管理过程，能够说明项目从任务分解到最终交付的演进路径。 & 仓库文档或交付材料目录。 \\
PE-22 & Markdown 过程资料 & “Markdown 过程资料”用于支撑软件工程管理过程，能够说明项目从任务分解到最终交付的演进路径。 & 仓库文档或交付材料目录。 \\
PE-23 & 外部数据调研 & “外部数据调研”用于支撑软件工程管理过程，能够说明项目从任务分解到最终交付的演进路径。 & 仓库文档或交付材料目录。 \\
PE-24 & 大模型接入说明 & “大模型接入说明”用于支撑软件工程管理过程，能够说明项目从任务分解到最终交付的演进路径。 & 仓库文档或交付材料目录。 \\
PE-25 & MCP 说明 & “MCP 说明”用于支撑软件工程管理过程，能够说明项目从任务分解到最终交付的演进路径。 & 仓库文档或交付材料目录。 \\
PE-26 & 最终 README & “最终 README”用于支撑软件工程管理过程，能够说明项目从任务分解到最终交付的演进路径。 & 仓库文档或交付材料目录。 \\
PE-27 & 风险登记册 & “风险登记册”用于支撑软件工程管理过程，能够说明项目从任务分解到最终交付的演进路径。 & 仓库文档或交付材料目录。 \\
PE-28 & 质量矩阵 & “质量矩阵”用于支撑软件工程管理过程，能够说明项目从任务分解到最终交付的演进路径。 & 仓库文档或交付材料目录。 \\
PE-29 & 变更控制 & “变更控制”用于支撑软件工程管理过程，能够说明项目从任务分解到最终交付的演进路径。 & 仓库文档或交付材料目录。 \\
PE-30 & 维护建议 & “维护建议”用于支撑软件工程管理过程，能够说明项目从任务分解到最终交付的演进路径。 & 仓库文档或交付材料目录。 \\
\bottomrule
\end{longtable}

\end{document}
